\documentclass[11pt]{article}

\usepackage[a4paper,margin=1in]{geometry}
\usepackage[utf8]{inputenc}
\usepackage[T1]{fontenc}
\usepackage{lmodern}
\usepackage{titlesec}
\usepackage{enumitem}
\usepackage{xcolor}
\usepackage{array}
\usepackage{longtable}
\usepackage{booktabs}
\usepackage{parskip}
\usepackage{fancyhdr}
\usepackage[colorlinks=true,linkcolor=blue!50!black,urlcolor=blue!50!black,citecolor=blue!50!black]{hyperref}

\definecolor{rulegray}{gray}{0.55}

\titleformat{\section}{\large\bfseries}{\thesection}{0.6em}{}
\titleformat{\subsection}{\normalsize\bfseries}{\thesubsection}{0.6em}{}
\titlespacing*{\section}{0pt}{1.4em}{0.6em}
\titlespacing*{\subsection}{0pt}{1.1em}{0.4em}

\pagestyle{fancy}
\setlength{\headheight}{14pt}
\fancyhf{}
\fancyhead[R]{\itshape Aashutosh Aryal}
\fancyfoot[C]{\thepage}
\renewcommand{\headrulewidth}{0.4pt}
\fancypagestyle{plain}{\fancyhf{}\fancyhead[R]{\itshape Aashutosh Aryal}\fancyfoot[C]{\thepage}\renewcommand{\headrulewidth}{0.4pt}}

\newcommand{\decision}[1]{\medskip\noindent\textbf{Decision: #1}\par\nopagebreak\smallskip}
\newcommand{\rationale}{\noindent\textbf{Rationale.}\ }
\newcommand{\tradeoff}{\noindent\textbf{Trade-off.}\ }

\begin{document}

% ---------------------------------------------------------------------
% Title page
% ---------------------------------------------------------------------
\begin{titlepage}
\centering
\vspace*{1.2cm}
{\sffamily\small\scshape\color{rulegray} Project Technical Report}\\[2.2cm]

{\Huge\bfseries git-feature}\\[0.6cm]
{\Large Storage, Modeling, and Derivation:\\[0.2cm] Design Decisions of a Feature-Oriented Git Extension}\\[2.2cm]

{\large Aashutosh Aryal}\\[0.4cm]

\vfill

{\large Version 1, August 2026}\\[0.3cm]
{\small\color{rulegray} feature-oriented-git}

\vspace*{1.5cm}
\end{titlepage}

% ---------------------------------------------------------------------
\section{Overview and Scope}

This report documents the technical decisions behind \emph{git-feature}, a
tool that lets code be labeled with the product features it implements,
keeps those labels attached correctly as history is rewritten, and can
derive a working build containing only a chosen combination of features.
The tool is a rewrite of \emph{Git with Features} (\emph{GWF}), an earlier
prototype that first explored recording \emph{feature associations} directly in Git.


The report is organized around the three parts of the system: the
\textbf{storage model} (Section~\ref{sec:storage}), which records which code
belongs to which feature and keeps that record accurate over time; the
\textbf{feature model} (Section~\ref{sec:modeling}), which describes what
features exist and which combinations of them are valid products; and
\textbf{variant derivation} (Section~\ref{sec:derivation}), which turns a
chosen combination into an actual, working codebase. Each section states a
decision, the reasoning behind it, and, where relevant, the trade-off it
accepts.

% ---------------------------------------------------------------------
\section{Rewriting Git with Features}
\label{sec:rewrite}

GWF established the core idea this project keeps: feature associations
should be recorded as an addition alongside a project's real history, not
by editing that history directly. Rewriting a commit to embed feature
metadata inside it would make ordinary history operations (reordering,
splitting, combining commits) destroy the very thing traceability is
supposed to survive. GWF calls this a \emph{non-invasive augmentation
strategy}: metadata is linked to commits through auxiliary structures,
without altering the commits themselves. Keeping that strategy was the
right starting point, and this project preserves it.

What motivated a rewrite rather than an extension of the original codebase
was a specific, structural problem in how GWF addressed that metadata. GWF
filed each feature association under the commit's content hash: a value
computed automatically from a commit's exact contents. That value is not a
stable name for a commit so much as a fingerprint of its current state,
and it changes the moment the commit is edited, reordered, or copied onto
another branch. An association filed under a hash that has since changed
does not raise an error; it simply stops matching anything, and the
association is silently lost. This is not a corner case: rebasing,
amending a commit message, and cherry-picking are all routine developer
actions, and all three break hash-keyed metadata.

Three further gaps shaped the scope of the rewrite:

\begin{itemize}[leftmargin=1.4em]
  \item GWF recorded, at most, a flat feature name attached to an entire
    commit. Its own evaluation names this directly: the prototype
    associates features at \emph{commit granularity}, and identifies
    finer-grained association as future work. A commit that touches two
    unrelated features had no way to express that split; the whole commit
    inherited a single label.
  \item There was no explicit model of what features exist, how they
    relate to each other, or which combinations form a legal product. A
    ``feature'' was just a string.
  \item There was no way to go from recorded associations to an actual
    build. GWF was a read-only traceability aid, not a build tool; its own
    discussion of related work draws exactly this line against tools that
    do support derivation, noting they operate largely outside version
    control instead.
  \end{itemize}

The rewrite keeps GWF's non-invasive augmentation strategy, replaces the
addressing scheme that broke under routine history rewrites, resolves the
commit-granularity limitation named above (Section~\ref{sec:regions}), and
adds the two capabilities (an explicit feature model, and derivation from
it) that did not exist before. These are documented in the sections that
follow.

% ---------------------------------------------------------------------
\section{The Storage Model}
\label{sec:storage}

\subsection{What Was Stored, Concretely}
\label{sec:storage-example}

Before going through each decision in turn, it helps to see one real
record from each system side by side, rather than only describing the
difference in the abstract. Both examples below come from running the
respective tool against the same small, two-commit project.

\medskip\noindent\textbf{GWF.} Tagging one commit with the feature
\texttt{auth} wrote a single file onto GWF's \emph{metadata branch}, an
ordinary Git branch named \texttt{feature-metadata}, at the path
\nolinkurl{auth/<commit-hash>/<timestamp>-<short-hash>}. Its content, a
small structured text format called JSON, was:

\begin{verbatim}
{
    "commit": "43932d3da78b667736f1af2d6a027dfed108b029",
    "authors": ["Dev"],
    "date": "2026-08-05T13:31:45.743465",
    "features": ["auth"],
    "changes": {
        "code_changes": [],
        "name_change": null,
        "constraint_changes": []
    }
}
\end{verbatim}

Two things about this record explain the problems described in
Section~\ref{sec:rewrite}. First, its location in the store, the folder
path itself, is built from the commit's content hash,
\texttt{43932d3d\ldots}: if that commit is later rebased, amended, or
cherry-picked, its hash changes, the folder this record lives in no
longer corresponds to anything, and the record is simply never found
again. Second, the record names exactly one feature for the whole commit; there
is no field that could describe ``these specific lines, not the rest of
the commit,'' the commit-granularity limitation GWF's own evaluation
names as future work.

\medskip\noindent\textbf{This project.} Tagging the equivalent commit with
the feature \texttt{core} writes two separate records to a ref outside the
branch namespace (Section~\ref{sec:branch-namespace}). The first maps the
commit's hash to a permanent identifier:

\begin{verbatim}
{
  "217f6ca05f87aaabaa03b944e2089bb92f5a2035": {
    "sha": "217f6ca05f87aaabaa03b944e2089bb92f5a2035",
    "change_id": "01KYPVNRPKW5BNQPEERY2PP86V",
    "patch_id": "7289c31690c6a128d4f54cfe274e56806f93214e"
  }
}
\end{verbatim}

The second is the annotation itself, filed under that identifier rather
than under the hash:

\begin{verbatim}
[
  {
    "anchor": {
      "change_id": "01KYPVPA4NRNNSEJA951YEEHM0",
      "path": "app.py",
      "old_span": [0, 0],
      "new_span": [1, 2],
      "fingerprint": "7fa8eae0365e947e"
    },
    "presence": "core",
    "author": "Dev",
    "ts": "2026-07-29T11:55:51.829972Z",
    "note": null
  }
]
\end{verbatim}

The difference is visible directly in the records. The annotation is
filed under \texttt{change\_id}, a value this tool assigns and controls,
not under the commit's hash; the small mapping file above is the only
place that ever needs to change if the commit is later rewritten, and the
automation in Section~\ref{sec:identity} updates exactly that file, and
nothing else, when it happens. The annotation also carries
\texttt{old\_span}/\texttt{new\_span}, which lines of that specific
historical change this label applies to, and \texttt{fingerprint}, a
content check used during resolution (Section~\ref{sec:resolution}), so
it describes one region of the commit rather than the commit as a whole.
\texttt{presence} holds a full condition (Section~\ref{sec:presence}), not
only a single feature name.

\subsection{A Metadata Store Outside the Branch Namespace}
\label{sec:branch-namespace}

\decision{Feature metadata lives in its own commit history, kept outside
Git's ordinary branch namespace, rather than on GWF's metadata branch.}

\rationale Git only treats a handful of naming conventions specially
(branches, tags, and a few others); nothing prevents a tool from writing
its own ordinary commit history under a name outside those conventions.
Doing so gets every property a branch would give (a durable, append-only,
content-addressed history, inspectable with the same tooling) without the
side effect of that history appearing in an everyday list of branches or
being checked out by mistake. GWF's metadata branch is a real, visible
branch for exactly this purpose; keeping the same non-invasive structure
while moving it out of the conventional namespace was a small change with
a direct benefit: nothing about the metadata store competes for the
developer's attention or can be confused with a branch of real work.

\tradeoff The store no longer benefits automatically from the tooling
that exists specifically for branches, most importantly the default rules
for what a plain push or fetch transfers. This is addressed directly in
Section~\ref{sec:sharing}.

\subsection{Stable Identity Under Rebase, Amend, and Cherry-Pick}
\label{sec:identity}

\decision{A commit's feature record is filed under a permanent,
tool-assigned identifier minted the first time the commit is seen, not
under the commit's content hash. Two small pieces of automation, run by
Git itself immediately after a commit is made and immediately after
history is rewritten, keep that identifier attached to the commit's
current form.}

\rationale This is the direct fix for the problem described in
Section~\ref{sec:rewrite}. Separating ``the identity a record is filed
under'' from ``the value Git derives from a commit's current content''
means the record can survive exactly the operations that discard the
latter. The identifier is minted lazily and cheaply: an unlabeled commit
gets one the first time it is either touched by the automation described
above, or explicitly annotated. Because Git already exposes hook points
that fire automatically after these specific actions, migrating the
identifier forward requires no discipline from the developer; it happens
as a side effect of the same commands they were already running.

A commit made somewhere this automation was not active, most commonly a
collaborator's clone before the tool was adopted there, or history that
predates adoption entirely, will not have an identifier yet. A
reconciliation step recovers one after the fact, first by comparing a
commit's exact set of changes against every already-known identifier's
changes (a reliable, unambiguous match when the change is unmodified),
and falling back to a text-similarity comparison when no exact match
exists. Genuinely ambiguous cases, several equally plausible matches, are
reported rather than resolved by guessing, and are left for a person to
confirm.

\tradeoff This identifier only has meaning within this project's
metadata store; it is not something Git itself, or any other tool, is
aware of. The reconciliation fallback is heuristic, not exact, for commits
the automation never saw.

\subsection{Regions in Place of Whole-Commit Associations}
\label{sec:regions}

\decision{A feature association is attached to the specific lines a
change introduced (a \emph{region}), not to the commit that introduced
them.}

\rationale Commits are a unit of convenience for the developer, not a
unit of feature scope. A single commit routinely touches more than one
concern: a bug fix bundled with an unrelated small refactor, or a change
that happens to serve two features at once. Attaching an association to
the whole commit forces every line in it to share one classification,
which is often simply false. Attaching it to the specific lines a change
introduced lets two unrelated edits landing in the same commit be
correctly and independently associated. GWF's finest granularity was the
commit; its own evaluation names this directly as a limitation and points
to finer-grained association as future work. A region is this project's
answer to that: the same concept of a feature association, GWF's term,
narrowed from a whole commit down to the lines that actually implement it.

A region is recorded relative to the historical diff that introduced it:
which file, which lines within that one diff, and a content fingerprint of
those lines. None of this needs to be updated later, because it describes
a historical fact and history does not change. What does change is
\emph{where that region lives now}, which is answered separately (see
Section~\ref{sec:resolution}).

\subsection{Presence Conditions}
\label{sec:presence}

\decision{A region's label is a boolean expression over feature names (a
\emph{presence condition}), for example \texttt{auth}, or
\texttt{payments \& premium \& !trial}, rather than a single feature
name.}

\rationale Code frequently only makes sense under a combination of
features, not one feature considered alone. Restricting every region to
exactly one feature name would force an artificial choice whenever a
region's real condition is a conjunction, a disjunction, or a negation of
several features. A presence condition is deliberately kept to a small
grammar, feature names, negation, conjunction, disjunction, and
parentheses, since expressing on/off combinations is the only thing this
expression needs to do; it does not need the structural vocabulary
(cardinalities, groups, required/excluded relationships) that the feature
model itself expresses (Section~\ref{sec:modeling}). Keeping the two
grammars separate, and the presence-condition grammar deliberately small,
keeps evaluating a region's fate under a chosen configuration a cheap,
unambiguous lookup.

\subsection{Resolving a Region at a Later Commit}
\label{sec:resolution}

\decision{Locating a region's current line range is computed on demand,
at the point it is needed, in a fixed sequence of increasingly
approximate strategies, each carrying an explicit confidence label rather
than a silent best guess.}

\rationale Line numbers are not a stable coordinate: every subsequent
edit to a file can shift, split, or reshape a previously recorded region.
Storing a line number once and trusting it forever would silently go
wrong the first time surrounding code changed. Instead, the current
location is derived fresh each time it is asked for, using the cheapest,
most reliable method first and falling back only when it fails:

\begin{enumerate}[leftmargin=1.6em]
  \item \textbf{Exact.} Trace each line of the target file back to the
    change that produced it; if the region's lines are still found
    together, in their originally recorded shape, the match is exact.
  \item \textbf{Moved.} The same trace finds the region's lines, but an
    unrelated edit has landed inside the region and split it into more
    than one contiguous piece. Every piece is still accounted for, just
    not as a single block.
  \item \textbf{Fuzzy.} If line-history tracing finds nothing, for
    instance after a heavy rewrite or reformat, fall back to searching
    file content directly for the region's recorded fingerprint, and
    failing that, for the closest text match above a fixed similarity
    threshold.
  \item \textbf{Lost.} None of the above succeeds. This is reported as
    what it is, not silently treated as ``nowhere'' or attached to
    whatever code happens to be nearby.
\end{enumerate}

\tradeoff Confidence can degrade over time as code is reshaped. This is
an honest limitation of any content-based tracking approach, not a defect
specific to this design; the alternative, silently trusting a stale
location, is strictly worse.

\subsection{A Worked Example}
\label{sec:resolution-example}

The four cases above are easiest to see concretely, one step at a time,
against a single small function tagged with the feature \texttt{bonus}:

\begin{verbatim}
def calc_bonus(rate):
    return rate * 100
\end{verbatim}

\begin{enumerate}[leftmargin=1.6em]
  \item \textbf{Right after tagging.} Nothing has changed yet. Asking
    where this region is now finds it immediately, unchanged, at its
    original position: \texttt{exact}.
  \item \textbf{A later, unrelated commit reformats the function,
    changing only whitespace:}
\begin{verbatim}
def   calc_bonus(rate) :
        return rate * 100
\end{verbatim}
    Tracing these lines back, the first and most reliable step, now points to the reformatting
    commit, not the one that introduced the feature, because every line
    in the region was rewritten, even though only its spacing changed.
    The trace comes back empty. But the region's content, with whitespace
    ignored, is exactly what was recorded when it was tagged, so a direct
    content search finds it again at the same position. This is reported
    as \texttt{fuzzy}, because it did not come from the reliable trace,
    even though the match itself is exact.
  \item \textbf{A later commit renames the parameter throughout the
    function:}
\begin{verbatim}
def calc_bonus(pct_rate):
    return pct_rate * 100
\end{verbatim}
    This defeats the content search too: with whitespace ignored, the
    text is genuinely different now, not just re-spaced. Comparing the
    originally recorded text to the candidate text at this position,
    however, still finds it about 90\% identical, comfortably over the
    fixed acceptance threshold (75\%), so this closest match is accepted.
    This is also reported as \texttt{fuzzy}; the label does not
    distinguish an exact content match from a close-enough one, only that
    neither came from the trace.
  \item \textbf{A later commit deletes the function and writes something
    unrelated in its place.} Nothing left in the file is even 75\%
    similar to the originally recorded text. Resolution reports
    \texttt{lost}, and, because this region needed to be removed under
    the configuration being built, derivation refuses outright rather
    than guess:
\begin{verbatim}
!! lost-region: region of 'bonus' in app.py must be removed
   but cannot be located at the base commit
1 conflict(s) -- variant cannot be materialized
\end{verbatim}
\end{enumerate}

This progression is also the practical answer to a specific failure mode
of trace-based tracking on its own: a single commit that reformats or
auto-formats a whole file would, under trace-based tracking alone,
reassign every line in it to that formatting commit, destroying every tag
those lines used to carry, even though nothing about the code's meaning
changed. Falling back to content and similarity search specifically does
not depend on the trace succeeding, so a pure reformat (case 2) is
recovered exactly and a moderate rewrite (case 3) is recovered
approximately; only a genuine replacement of the underlying content (case
4) is ever reported as lost, and even then, honestly, instead of guessing.

\subsection{Sharing and Merging the Store}
\label{sec:sharing}

\decision{Transferring the metadata store between clones, and combining
two independently modified copies of it, are explicit operations with
their own deterministic merge rules, rather than something a plain push
or fetch handles automatically.}

\rationale This is the direct consequence of Section~\ref{sec:storage}'s
first decision. Git's default transfer rules only mirror the conventional
branch namespace; moving the store outside that namespace was necessary to
keep it out of the way, but it also means nothing built into Git already
knows to carry it along. The alternative, keeping the store on an
ordinary branch specifically so it inherits default transfer behavior,
was rejected because it reintroduces the exact clutter and
accidental-checkout risk the first decision in this section was made to
avoid. Given that trade-off, the store's synchronization is instead
implemented directly: explicit commands move it between clones, and a
merge routine combines two diverged copies using rules chosen so that two
independent machines always converge on the same result (for example: an
identifier minted twice, once on each side, for the same commit keeps a
single, deterministic choice between the two, and every downstream record
that referenced the discarded one is transparently redirected).

\tradeoff Sharing metadata is a step a developer has to remember to take;
it is not a free side effect of sharing code the normal way.

\subsection{Known Limitations of the Storage Model}
\label{sec:storage-limits}

\begin{itemize}[leftmargin=1.4em]
  \item The identity-tracking automation described in
    Section~\ref{sec:identity} is local to whichever clone it is
    installed on. Commits made elsewhere depend on the reconciliation
    fallback, which is reliable for unmodified changes but can leave
    genuinely ambiguous cases for a person to resolve.
    In practice this means fetching and adopting an existing store (Section~\ref{sec:sharing}) before initializing or backfilling a fresh clone, since that sidesteps reconciliation entirely for anything the store already knows about; only commits no clone's hooks ever touched at all still require it
  \item Resolution confidence (Section~\ref{sec:resolution}) is
    best-effort and can degrade with sufficiently heavy rewrites; this is
    surfaced honestly rather than hidden, but it means traceability
    quality is not guaranteed to remain perfect indefinitely.
  \item Merging two copies of the store (Section~\ref{sec:sharing}) is
    only automatic when both were written by compatible versions of the
    tool's internal record format; an incompatible pair fails closed
    rather than risk a misinterpreted merge.
  \item Nothing in the storage model infers a region's feature
    automatically. A change that is never labeled is permanently treated
    as unconditional code; this is a deliberate default, not an
    oversight, but it does mean coverage is only as complete as what a
    developer actually records.
  \item The similarity fallback within resolution (step 3 of
    Section~\ref{sec:resolution}) accepts the single best match above its
    threshold without checking whether a second, comparably good match
    exists elsewhere in the project. This is unlike the ambiguity check
    performed during identifier reconciliation (Section~\ref{sec:identity}),
    which explicitly reports a tie for a person to resolve; a resolution
    that happens to match more than one plausible location is, at
    present, decided silently rather than flagged.
  \item A presence condition's syntax (Section~\ref{sec:presence}) is not
    checked at the moment it is recorded. A malformed condition is only
    caught later, when the feature model is explicitly validated or a
    variant that depends on it is derived.
\end{itemize}

% ---------------------------------------------------------------------
\section{The Feature Model}
\label{sec:modeling}

\subsection{Placing the Model File at the Project Root}
\label{sec:model-root}

\decision{The feature model is a single file, \texttt{model.cfr}, kept at
the top level of the project's working directory, tracked and versioned
exactly like any other source file.}

\rationale Two alternatives were rejected:

\begin{itemize}[leftmargin=1.4em]
  \item Keeping the model in the metadata store (Section~\ref{sec:storage})
    alongside everything else would hide a decision that is, by nature,
    about the product as a whole from anyone simply browsing the
    repository, with no indication it exists at all unless they already
    knew where to look.
  \item Nesting it inside a configuration subdirectory would demote it to
    the same tier as incidental tooling configuration, when it is closer
    in kind to a project's top-level manifest, the same reason a
    \texttt{README} or package manifest sits at the root.
  \end{itemize}


Placing it at the root instead means it is the first thing visible in a
directory listing, requires no special handling to version, diff, or
review (it participates in ordinary code review the same as any other
file change), and is trivially available to the derivation pipeline
without any extra bookkeeping: resolving a configuration simply reads
\texttt{model.cfr} out of the commit tree being built from
(Section~\ref{sec:pipeline}), the same way it would read any other
tracked file.

\tradeoff A root-level file claims a name that could, in principle,
collide with an unrelated project convention. This was judged an
acceptable, low-probability cost against the visibility gained.

\subsection{A Deliberately Small Subset of Clafer}

\decision{The feature model is written in a small, explicitly bounded
subset of the Clafer modeling language: a feature tree with
cardinalities, group cardinalities, boolean cross-tree constraints, and
named, concrete configurations. Everything else Clafer supports
(references, arithmetic, quantifiers, structural joins, enumerations) is
deliberately excluded and fails to parse.}

\rationale Clafer itself was chosen over inventing a bespoke format
because it is an established, purpose-built notation for exactly this
problem, feature trees, group relationships, and constraints, with a
syntax concise enough to sit comfortably in a small text file. Adopting
only a subset of it, rather than the full language, was a separate and
equally deliberate choice: nothing beyond ``which features exist, how they
relate, and which combinations are valid'' is needed to annotate code and
derive variants, and a smaller grammar is easier to implement correctly,
easier to keep correct as the tool evolves, and easier for a modeler to
learn completely rather than partially. The subset restriction is stated
directly in the grammar itself, so an author who reaches for a feature
outside it is told immediately, at parse time, rather than discovering
a silent misinterpretation later.

\subsection{Resolving Named Configurations}

\decision{Resolving a configuration works out every feature's status, not
just the ones it names directly, by repeatedly applying the model's
rules until nothing new follows. If the rules ever contradict each
other, the resolver rejects the configuration and says exactly why.}

\rationale A configuration like \texttt{Full} only names a few features
directly. The rest follow from rules already in the model, such as
``this feature requires that one'' or ``exactly one of these three must
be on.'' One pass through these rules is often not enough: settling one
feature can complete a group, or satisfy a rule, that then settles
another. So the resolver keeps re-checking all the rules until a full
pass changes nothing, and only then treats anything still undecided as
excluded. If two rules ever force the same feature to different values,
that is a genuine contradiction in the model, and the resolver stops and
reports exactly what conflicts, instead of guessing.


% ---------------------------------------------------------------------
\section{Variant Derivation}
\label{sec:derivation}

\subsection{Design Principle: Derive, Never Maintain}

\decision{A variant (a build containing only a chosen configuration's
features) is computed fresh from history every time it is requested. It
is never kept as a long-lived branch that is separately maintained and
periodically updated.}

\rationale A permanently maintained variant branch drifts the moment
anything upstream changes, and keeping it current becomes a manual
responsibility someone has to remember to discharge, precisely the kind
of bookkeeping this project exists to remove. Recomputing a variant
on demand from the tagged history, instead, means it can never be stale:
it always reflects exactly what the current (or a chosen historical)
state of the tagged history says, with no separate copy to fall out of
sync.

\subsection{The Four-Stage Pipeline}
\label{sec:pipeline}

\decision{Deriving a variant is a pipeline of four distinct, individually
checkable stages, rather than a single opaque step.}

\begin{enumerate}[leftmargin=1.6em]
  \item \textbf{Configure.} The requested configuration is resolved
    against \texttt{model.cfr} as it exists in the tree of the commit
    being built \emph{from}, not the working directory, so that
    deriving from an older point in history uses the model as it existed
    at that point. This yields a total feature-to-boolean assignment,
    or an outright rejection if the configuration is illegal.
    \item \textbf{Project.} For every region belonging to that history,
    decide whether it stays or goes by checking its presence condition
    against the assignment, and find where it currently lives
    (Section~\ref{sec:resolution}). These decisions are then checked
    against each other, and any conflict is named rather than silently
    resolved: two regions disagreeing over the same lines, a kept region
    sitting inside one about to be removed, or a region that must be
    removed but cannot be found.
  \item \textbf{Materialize.} The reduced tree is built: every disabled
    region's lines are spliced out of the files that contain them, and
    every file untouched by any removal is carried over unchanged. This
    produces an actual, ordinary branch containing the result.
  \item \textbf{Verify.} Before that branch becomes the visible result,
    it gets a second, independent check: for each region, its actual
    recorded content, not the resolver's answer about where it should
    be, is looked for in the produced tree, confirming a kept region is
    really there and a removed one is really gone. Any mismatch aborts
    the operation outright.
  \end{enumerate}

\rationale Splitting derivation into stages with individually checkable
outputs, rather than one combined operation, means each kind of failure
(an illegal configuration, a contradictory set of regions, a
mis-resolved location) is caught at the earliest point it can be
detected, with a specific, actionable reason, instead of surfacing later
as an unexplained bad build. A preview mode that stops after the second
stage and reports the plan without writing anything follows directly from
this separation: previewing a derivation is simply not running the last
two stages.

\subsection{Verification as a Gate, Not a Formality}

\rationale The materialize stage trusts the resolver's answer about
where a region currently lives; that resolver, as described in
Section~\ref{sec:resolution}, is explicitly best-effort past its first
tier. Treating its output as automatically correct would let a
low-confidence resolution silently produce a subtly wrong build. The
verify stage exists specifically to not extend that trust: it recovers
what a region's content actually was, independently of the location the
resolver reported, and checks that content's presence or absence in the
result, so a derivation is only ever handed over once it has been
confirmed to actually contain what it claims to.

\tradeoff This check operates purely on text content; it confirms the
expected lines are present or absent, not that the resulting code still
compiles or behaves correctly. It catches a specific, important class of
failure (a wrong or incomplete removal), not every possible one.

\subsection{Editing a Derivation: View and Putback}

\decision{A derived variant can optionally be opened as an editable
session (a \emph{view}), and an edit made to it can be mapped back
automatically onto the corresponding position in the full, undereived
history (\emph{putback}), rather than requiring a plain, disposable
derivation only.}

\rationale A reduced variant is not only useful as a shipping artifact;
it is also, simply by being smaller, sometimes an easier place to notice
and fix a bug than the full codebase would be. Without a way to carry a
fix made there back to the real history, that advantage would be wasted,
since a plain derivation is disposable and rebuilding it would discard
any edit made directly on it. A view therefore additionally records
exactly which line ranges of the source were removed to produce it; that
record is precisely the information needed to translate a position in the
reduced file back into a position in the full file, and vice versa.
Applying a mapped edit back is then an ordinary commit on the real
history, immediately eligible for its own feature labels like any other
commit.

\tradeoff Only edits that map unambiguously are accepted: an insertion
or an in-place change that sits entirely within one surviving piece of a
file. An edit that straddles the boundary between kept and removed code,
or a change that is not a simple content edit (deleting a file, a binary
change), is refused rather than partially applied, and the mapping
itself becomes stale, and must be refreshed, if the real history moves on
while a view is open.

\subsection{Known Limitations of Derivation}

\begin{itemize}[leftmargin=1.4em]
  \item Region removal operates purely on text; the pipeline has no
    understanding of the target language's syntax or semantics, so a
    poorly chosen region boundary can still produce output that passes
    content verification (Section~\ref{sec:pipeline}) yet fails to
    compile: removing a helper function while a call to it survives in a
    kept region leaves exactly the text each region promised, so
    verification passes, even though the result no longer runs.
  \item Contradictions found during the project stage are refused, not
    resolved automatically; an automatic resolution could easily choose
    wrong, so some configurations require a person to intervene before a
    build can be produced at all.
  \item Putback (Section~\ref{sec:derivation}) only accepts straightforward
    edits, as described above; more complex edits are rejected outright
    rather than approximated.
\end{itemize}

% ---------------------------------------------------------------------
\section{Summary of Key Concepts}

Table~\ref{tab:concepts} collects the terms introduced in this report as
a quick reference.

\begin{longtable}{@{}p{0.24\textwidth}p{0.68\textwidth}@{}}
\toprule
\textbf{Term} & \textbf{Meaning} \\
\midrule
\endhead
Metadata store & The separate, non-branch commit history holding all
  feature records (Section~\ref{sec:storage}). \\
Permanent identifier & The tool-assigned, hash-independent identity a
  commit's records are filed under (Section~\ref{sec:identity}). \\
Region & A specific line range from one historical change, labeled with a
  presence condition (Section~\ref{sec:presence}). \\
Presence condition & A boolean expression over feature names, attached to
  a region, deciding when it should be present (Section~\ref{sec:presence}). \\
Resolution confidence & How certain the tool is about a region's current
  location: exact, moved, fuzzy, or lost (Section~\ref{sec:resolution}). \\
Configuration & A total, resolved feature-to-boolean assignment produced
  from a named instance in the model (Section~\ref{sec:modeling}). \\
Variant & A derived, reduced build containing only the regions a chosen
  configuration keeps (Section~\ref{sec:derivation}). \\
Projection & The stage that decides every region's fate under a
  configuration and reports contradictions (Section~\ref{sec:pipeline}). \\
View / putback & An editable derivation, and the mechanism mapping an
  edit made on it back onto the full history. \\
\bottomrule
\caption{Key terms used throughout this report.}
\label{tab:concepts}
\end{longtable}

\section{Summary of Design Decisions}

\begin{longtable}{@{}p{0.32\textwidth}p{0.6\textwidth}@{}}
\toprule
\textbf{Decision} & \textbf{Reason, in short} \\
\midrule
\endhead
Metadata store outside the branch namespace & Keeps records durable and
  inspectable without cluttering everyday Git commands. \\
Records filed under a permanent identifier & Survives rebase, amend, and
  cherry-pick, which broke hash-keyed records. \\
Region-level, not commit-level, labels & A commit can serve more than one
  feature; a label should not force one classification on all of it. \\
Presence conditions as boolean expressions & Code often depends on a
  combination of features, not one in isolation. \\
Tiered, confidence-labeled resolution & Line positions are not stable;
  honesty about uncertainty beats a silent guess. \\
Explicit store sharing and merging & The cost of moving off the branch
  namespace, addressed directly rather than left unsolved. \\
\texttt{model.cfr} at the project root & Visible by default, versioned
  like any source file, no special tooling required to find or diff it. \\
A small, bounded Clafer subset & Only what feature modeling actually
  needs; smaller surface, easier to implement and learn completely. \\
Repeatedly apply the model's rules to resolve a configuration & Settles
  features not named directly, and rejects impossible configurations by
  name instead of guessing. \\
Derive on demand, never maintain a copy & A recomputed variant cannot go
  stale the way a hand-maintained one would. \\
Four-stage pipeline with a verification gate & Isolates each kind of
  failure and refuses to hand over an unverified result. \\
View and putback & Lets a fix found in a smaller, reduced build return to
  the real history automatically. \\
\bottomrule
\caption{Design decisions covered in this report, condensed.}
\end{longtable}

\end{document}
